% To compile using latexmk on the command line, run the following: 
% latexmk -pdf outline.tex

\documentclass[11pt]{article}

% ---------- Page formatting ----------
\usepackage{geometry}
\geometry{margin=1in}

\usepackage{lmodern}
\usepackage[T1]{fontenc}
\usepackage{setspace}
\onehalfspacing

% ---------- Graphics ----------
\usepackage{graphicx}
\usepackage{float}

% ---------- Math ----------
\usepackage{amsmath}
\usepackage{amssymb}
\usepackage{bm}

% ---------- ACS citations ----------
\usepackage{achemso}
% \setkeys{acs}{articletitle=true}

% ---------- Title ----------
\title{\textbf{Computational Identification of High-Altitude Molecules Using Molecular Networking and Liquid Chromatography High-Resolution Mass Spectrometry (LC-HRMS): Developing Software by Chemists, for Chemists}}
\author{Raelyn Brooks}
\date{\today}

\begin{document}

\maketitle

\nocite{*}

\begin{doublespace}

\section{Abstract}
    Written after the results are complete. Summarize the environmental problem, research question, methods, software contribution, computational evaluation, chemical case studies, principal findings, and limitations.

\section{Chapter 1: Introduction and Theoretical Background}
    \subsection{Research Problem and Motivation}
        \begin{itemize}
            \item Organic contaminants in remote, high-altitude environments
            \item The challenge of identifying unknown molecular features
            \item The need for reliable and reproducible molecular networking
        \end{itemize}

    \subsection{Environmental and Chemical Background}
        \begin{itemize}
            \item Long-range transport of organic contaminants
            \item High-altitude and remote surface-water environments
            \item Previous Wooster research and the current identification gap
        \end{itemize}

    \subsection{LC-HRMS, Suspect Screening, and Non-Targeted Analysis}
        \begin{itemize}
            \item LC-HRMS and tandem mass spectrometry
            \item Molecular features versus identified compounds
            \item Suspect screening and non-targeted analysis
            \item Limitations of current identification approaches
        \end{itemize}

    \subsection{Molecular Networking and Theoretical Foundations}
        \begin{itemize}
            \item MS/MS fragmentation and spectral similarity
            \item Cosine similarity and molecular-network construction
            \item Graph theory: nodes, edges, and connected components
            \item Network visualization and interpretation
            \item Limitations of spectral similarity as evidence of chemical identity
        \end{itemize}
    
    \subsection{Existing Software and Research Gap}
        \begin{itemize}
            \item MS-DIAL and GNPS in the existing workflow
            \item Documented limitations in the Wooster research workflow
            \item The need for software designed around chemists' needs
        \end{itemize}
    
    \subsection{Research Question and Specific Aims}
        Central research question: What computational workflow and interface are needed to make molecular networking reliable, reproducible, and useful for Wooster's environmental HRMS research?
        \begin{itemize}
        \item Aim 1: Define the chemistry and data requirements.
        \item Aim 2: Develop an end-to-end molecular-networking workflow.
        \item Aim 3: Evaluate computational correctness, reproducibility, performance, and workflow usefulness.
        \item Aim 4: Apply the software to environmental HRMS case studies.
        \end{itemize}
    
    \subsection{Project Scope and Thesis Contributions}
        \begin{itemize}
            \item Clarify what the software will and will not do.
            \item Identify the expected Chemistry and Computer Science contributions.
            \item Explain the organization of the thesis.
        \end{itemize}

\section{Chapter 2: Methodology, Research Design, and Software Development}
    \subsection{Overall Research Design and Workflow}
        \begin{itemize}
            \item Overall research approach and data flow
            \item Relationship between the chemistry and computational components
            \item Research datasets and evaluation strategy
        \end{itemize}

    \subsection{Environmental Samples and Chemistry Experimental Methods}
        \begin{itemize}
        \item Sample origins, locations, and sample types
            \item Instrumentation and analytical conditions
            \item LC-HRMS and MS/MS acquisition
            \item Positive- and negative-ion modes, where applicable
            \item Clarification of work performed by previous researchers 
        \end{itemize}
    
    \subsection{Upstream Data Processing and Chemical Annotation}
        \begin{itemize}
            \item MS-DIAL processing and version
            \item Feature detection, alignment, and spectral exports, as applicable
            \item Data-quality checks and inclusion criteria
            \item Suspect screening and candidate annotations
            \item Input data preparation and metadata
            \item Identification confidence and limitations
        \end{itemize}

    \subsection{Software Requirements and System Architecture}
        \begin{itemize}
            \item Functional and nonfunctional requirements
            \item Supported input formats and validation
            \item System modules and data flow
            \item Data structures and metadata representation
            \item Reproducibility and provenance
        \end{itemize}

    \subsection{Spectral Comparison and Similarity Calculation}
        \begin{itemize}
            \item Spectrum representation and preprocessing
            \item Peak matching and mass tolerances
            \item Cosine similarity implementation
            \item Edge thresholds and computational considerations
            \item Handling missing, incomplete, or unusable spectra
        \end{itemize}

    \subsection{Molecular-Network Construction}
        \begin{itemize}
            \item Node and edge representation
            \item Similarity thresholds and graph filtering
            \item Connected components and isolated features
            \item Linking network elements to chemical metadata
        \end{itemize}

    \subsection{User Interface Design and Interaction}
        \begin{itemize}
            \item Network layout and visual encodings
            \item Node and edge selection
            \item Feature metadata and similarity inspection
            \item Filtering, navigation, and network exploration
            \item Saving, exporting, or persistence, where implemented
            \item Design decisions informed by chemist workflows
        \end{itemize}

    \subsection{Software Implementation and Reproducibility}
        \begin{itemize}
            \item Programming languages, libraries, and dependencies
            \item Modular design, version control, and documentation
            \item Error handling and logging
            \item Configuration, parameter recording, and reproducible execution
        \end{itemize}
    
    \subsection{Software Evaluation and Testing Methods}
        \begin{itemize}
            \item Unit, integration, and end-to-end tests
            \item Reference calculations for similarity validation
            \item Graph-construction correctness
            \item Repeated-run reproducibility
            \item Runtime and scalability measurements
            \item Interface and workflow task evaluation
        \end{itemize}

    \subsection{Chemical Case-Study Design (Testing the Software on Real Data)}
        \begin{itemize}
            \item Selection of representative environmental samples
            \item Criteria for selecting samples, features, and network clusters
            \item Use of existing annotations and spectral evidence
            \item Approach to interpreting known and unidentified features
            \item Criteria for limiting chemical claims
        \end{itemize}

\section{Chapter 3: Results and Data Analysis}
    \subsection{Environmental Dataset and Feature Summary}
        \begin{itemize}
            \item Dataset characteristics and sample coverage
            \item Number of input features and usable MS/MS spectra
            \item Available annotations, exclusions, and data limitations
        \end{itemize}
    
    \subsection{Software Implementation Results}
        \begin{itemize}
            \item Completed workflow and supported functionality
            \item Implemented architecture and interface
            \item Representative input-to-network examples
        \end{itemize}

    \subsection{Computational Correctness and Validation}
        \begin{itemize}
            \item Similarity calculations versus reference values
            \item Input-validation and error-handling results
            \item Network-construction tests
            \item Reproducibility and persistence results
        \end{itemize}

    \subsection{Software Performance and Workflow Evaluation}
        \begin{itemize}
            \item Runtime and memory use, where measured
            \item Performance across tested dataset sizes
            \item Interface and workflow task results
            \item Requirements met, partially met, or unmet
        \end{itemize}

    \subsection{Molecular-Network Structure and Characteristics}
        \begin{itemize}
            \item Nodes, edges, connected components, and isolated features
            \item Similarity-score distributions
            \item Effects of thresholds and filtering, where evaluated
        \end{itemize}

    \subsection{Chemical Case Study 1: Annotated Molecular Features}
        \begin{itemize}
            \item Selected network and chemical context
            \item Spectral relationships and available annotations
            \item Interpretation and evidence strength
        \end{itemize}
    
    \subsection{Chemical Case Study 2: Unidentified Molecular Features}
        \begin{itemize}
            \item Selected unknown features or network clusters
            \item Relationships to annotated or other unidentified features
            \item Potential chemical relationships and unresolved questions
        \end{itemize}

    \subsection{Summary of Principal Results}
        \begin{itemize}
            \item Concise synthesis of the strongest computational and chemical findings
        \end{itemize}

\section{Chapter 4: Discussion}
    \subsection{Interpretation and Integration of Results}
        \begin{itemize}
            \item Answer the central research question using both result sets.
            \item Explain how chemistry informed software design and evaluation.
            \item Explain how the software supported chemical investigation.
        \end{itemize}
    
    \subsection{Computational Contributions and Design Trade-Offs}
        \begin{itemize}
            \item Reliability, correctness, reproducibility, and usability
            \item Strengths and weaknesses of algorithms and architecture
            \item Comparison with intended requirements and existing workflows
        \end{itemize}
    
    \subsection{Chemical Significance and Comparison with Literature}
        \begin{itemize}
            \item Interpretation of case-study findings
            \item Relationship to existing environmental chemistry research
            \item What networking reveals about unknown features
            \item Limits on claims about identity and transformation products
        \end{itemize}
    
    \subsection{Study and Implementation Limitations}
        \begin{itemize}
            \item Analytical and data-quality limitations
            \item Computational and scalability limitations
            \item Visualization and workflow limitations
            \item Evaluation constraints and limits on generalization
        \end{itemize}
    
    \subsection{Implications for Future Environmental HRMS Research}
        \begin{itemize}
            \item How the workflow may support future Wooster projects
            \item Additional data, validation, or software capabilities needed
        \end{itemize}    
    
\section{Chapter 5: Conclusions and Future Work}
    \subsection{Summary of the Research}
        Revisit the environmental problem, objectives, and approach.
    \subsection{Principal Computer Science Contributions}
        Summarize the implemented workflow and evaluation findings.
    \subsection{Principal Chemistry Contributions}
        Summarize the supported chemical findings and their significance.
    \subsection{Answer to the Central Research Question}
        State what the project demonstrated and what remains unresolved.
    \subsection{Future Work/Directions}
        Additional validation, environmental datasets, software improvements, and further chemical investigation.
    \subsection{Final Remarks}
        Conclude with the project's contribution to computational chemistry.

\end{doublespace}

% ---------- References ----------
\newpage
\bibliography{references.bib}

% ---------- Appendices ----------
\appendix

\section{Supported Input Formats and Data Schema}
\section{Additional Computational Methods and Test Cases}
\section{Supplementary Software Evaluation Results}
\section{Additional Molecular-Network Figures and Chemical Data}
\section{Software Setup and Reproducibility Instructions}
\section{AI Use Documentation, as Required by Departmental Policy}

\end{document}
